% TOP_PROBLEM: {"id": 1100213, "title": "Questions in Geometric Group Theory — Q 2.13", "category_id": 17, "category": {"id": 17, "name": "group_theory", "display_name": "Group Theory", "description": "Problems about groups, group actions, representations, and related algebraic structures.", "slug": "group-theory", "order_index": 17, "created_at": "2026-07-31T15:26:25.671Z"}, "status": "open", "problem_number": "AMR-010-0213", "set_id": 13}
% FIRST_POSTED: 2026-10-01
% ENTRY_KIND: statement_only
% UnsolvedMath ID 1100213; source catalogue code AMR-010-0213
% !TeX program = lualatex
% Definition and sources only; this file is not an LLM solution attempt.
\documentclass[11pt,a4paper]{article}
\usepackage[margin=25mm]{geometry}
\usepackage{fontspec}
\setmainfont{lmroman10-regular.otf}[BoldFont=lmroman10-bold.otf,
 ItalicFont=lmroman10-italic.otf,BoldItalicFont=lmroman10-bolditalic.otf]
\usepackage{amsmath,amssymb,mathtools}
\usepackage{unicode-math}
\setmathfont{latinmodern-math.otf}
\AtBeginDocument{\let\setminus\smallsetminus}
\usepackage{xurl,hyperref}
\hypersetup{colorlinks=true,urlcolor=blue}
\setcounter{secnumdepth}{0}
\setlength{\parindent}{0pt}
\setlength{\parskip}{5pt}
\setlength{\emergencystretch}{3em}
\begin{document}
\section*{Questions in Geometric Group Theory — Q 2.13}\label{1100213}
{\small UnsolvedMath ID 1100213 · AMR-010-0213 · Group Theory · AMR Open Problem Lists}
\subsection{Definitions and mathematical statement}
A Coxeter matrix on a finite set $S$ is a symmetric matrix with $m_{ss}=1$ and $m_{st}\in\{2,3,\ldots\}\cup\{\infty\}$ for distinct indices. Its group is
\[
 W(M)=\langle S\mid s^2=1\ (s\in S),\ (st)^{m_{st}}=1
                  \ (s\ne t,\ m_{st}<\infty)\rangle.
\]
The symbol $\infty$ records the absence of a relation for that pair. The finite labeled diagram, or equivalently the matrix, is the input describing a Coxeter presentation.

Classify these groups up to abstract isomorphism. More precisely, give necessary and sufficient conditions on two finite Coxeter matrices $M,N$ for $W(M)\cong W(N)$, in terms of their defining data or intrinsic invariants that can be determined from those data. An effective classification would in particular decide isomorphism from two finite matrices.

The requested isomorphism is not required to send the displayed generators to displayed generators, or even to reflections of the second Coxeter system. A reflection relative to $S$ means a conjugate of an element of $S$. Hence a classification only of marked Coxeter systems, or only of reflection-preserving isomorphisms, is a restricted version. Relabeling an identical diagram certainly produces an isomorphic group, but the question asks for all possible coincidences between presentations. Both finite and infinite Coxeter groups, including reducible diagrams, belong to the class.
\subsection{Short English statement}
Give a complete criterion for when two finite Coxeter presentations define isomorphic abstract groups.
\subsection{Sources}
\begin{itemize}
\item[S1] ULAM.AI and the UnsolvedMath Contributors, supplied problems.json, record 1100213 (AMR-010-0213), AMR Open Problem Lists. Catalogue identity and source transcription; CC BY 4.0.
\url{https://huggingface.co/datasets/ulamai/UnsolvedMath}.
\item[S2] Bestvina - Questions in Geometric Group Theory (pdf) (2004), Question 2.13 (PDF page 9). Original source indexed by AMR-010-0213.
\url{https://www.math.utah.edu/~bestvina/eprints/questions-updated.pdf}.
\end{itemize}
\end{document}
